% Transcription of folder 149, batch 01 (archivists' pages 1-20).
% Pass: Opus 5.5 (claude-opus-5-5), 2026-10-03 — first pass, unchecked against the pages by a human.
%
% Page 1 carries only a pencilled label ("Pages 1-67") and no mathematics; it is
% absent from the transcription. Grothendieck numbers every other sheet
% (author's pp. 1-10 = archivists' pp. 2, 4, ..., 20); the unnumbered sheets in
% between continue the text without a break.
\documentclass[11pt,a4paper]{article}
\input{../preamble/grothendieck}

\folder{149}
\batch{1}
\pages{1}{20}
\foldertitle{[Autour de La "Longue Marche" à travers la théorie de Galois, pages 1 à 67] : notes manuscrites (s.d.).}
\dating{s.d. — le groupe « Autour de La "Longue Marche" à travers la théorie de Galois » (141 à 149) est daté [à partir de 1978]-1983}
\watermark{Édition de démonstration}

\begin{document}

\section{1. Résultats de fidélité}

\page{2}\note{p. 1 de l'auteur}
À tout corps $K$, associons son topos étale $B_K$, qui est un topos (pro\add{fini}) galoisien. Le groupoïde des points de $B_K$ est noté $\Pi_K$, il est antiéquivalent canoniquement à la catégorie des clôtures alg. séparables de $K$. Si $\overline{K}$ est une telle clôture, son groupe des \add{$K$-}automorphismes $\operatorname{Gal}(\overline{K}/K)$ ou $E_{\overline{K}/K}$ s'identifie au groupe des \struck{$K$-}automorphismes du pt \add{de $B_K$} associé à $\overline{K}/K$ (il vaut peut-être mieux de dire : à l'opposé de ce groupe -- \struck{$\overline{K}$} la variance des clôtures alg. de $K$ est comme celle des foncteurs fibres, i.e. l'opposée de celle des points…)

Bien entendu, $B_K$ se reconstitue : à partir de $\Pi_K$, comme le topos des systèmes locaux (continus) sur $\Pi_K$ -- et en termes de $E_{\overline{K}/K}$, comme le topos des \uncertain{ensembles} discrets à action continue de $E_{\overline{K}/K}$.

Pour un hom de corps $K \to K'$,

\page{3}
i.e. un hom de schémas $\operatorname{Spec} K' \to \operatorname{Spec} K$, on a un morphisme de topos correspondant
\[
  (1)\qquad B_{K'} \longrightarrow B_{K^{*}}
\]
associé à un hom de groupoïdes fond.
\[
  (\uncertain{2})\qquad \Pi_{K'} \longrightarrow \Pi_{K^{*}} .
\]
\note{les deux numéros se ressemblent sur la page ; le second est lu « (2) » sous réserve. L'astérisque en exposant de $K$ est une petite marque que nous rendons par $K^{*}$ ; elle semble ne désigner que $K$ lui-même.}
Ceci s'explicite en disant qu'un \struck{clôture} objet de $\Pi_{K'}$ (i.e. pt. de $B_{K'}$, ou rev. universel de $B_{K'}$, ou clôture sép. \add{$\overline{K}'$} de $K'$) en définit un de $\Pi_{K^{*}}$ (ainsi, on prend $\overline{K}^{*}$ = clôture \add{alg.} \struck{sép.} de $K$ dans $\overline{K}'$) et pour deux pts correspondants, on a un hom de groupes fondamentaux correspondants, qui s'interprète p.ex. comme
\[
  (3)\qquad E_{\overline{K}'/K'} \longrightarrow E_{\overline{K}^{*}/K^{*}}
\]
et qui permet de reconstituer l'hom. de topos comme une « restriction des scalaires ».

L'image de (3) est le sous-groupe fermé de $E_{\overline{K}/K}$ qui correspond à la sous-extension

\page{4}\note{p. 2 de l'auteur}

\begin{tikzcd}[column sep=small, row sep=small]
K_1 \arrow[r, no head] \arrow[dr, no head] & K' \arrow[dr, no head] & \\
 & \overline{K} \arrow[r, no head] & \overline{K}'
\end{tikzcd}

$K_1$ de $\overline{K}/K$, clôture \add{alg.} \struck{sép.} de $K$ dans $K'$, i.e. $K_1 = \overline{K} \cap K'$. Quand $K'$ est une extension de type fini de $K$, $K_1$ est une ext. finie de $K$, et on en conclut que l'image de (3) est alors un sous-groupe d'indice fini, égal à $E_{\overline{K}/K}$ sss \add{$K_1 = K$ i.e.} $K$ est séparablement alg. clos dans $K'$. D'ailleurs, on montre sans mal que \add{(si $K'$ est ext. de type fini)} l'hom (3) est injectif sss $K'$ est une extension algébrique de $K$. Donc il est bijectif sss $K'$ est une extension radicielle de $K$. Dans la suite nous nous bornons (provisoirement) aux corps de car. 0, et la condition précédente signifie \add{alors} que $K \to K'$ est un \emph{isomorphisme}.

Ainsi, le foncteur $K \mapsto B_K$ ou $K \mapsto \Pi_K$, ou $(K, \overline{K}) \mapsto E_{\overline{K}/K}$, est \emph{conservatif} quand on se limite aux morphismes de corps \add{$K \to K'$ (de car. 0)} à ceux pour lesquels $K'$ est une ext. de type fini de $K$,

\page{5}
P.ex. il suffit de se limiter aux extensions de type fini du corps premier $\mathbb{Q}$ -- on trouve un foncteur conservatif de la cat. de ces corps, vers celle des groupoïdes \add{(ou des topos)}, au sens : un morphisme de corps qui donne une équivalence de groupoïdes (ou de topos) est un \emph{iso}.

\marginal{\ill{} est conservatif \ill{} \uncertain{catégorie} quotient de celle des corps \uncertain{où} une extension radicielle finie est considérée comme triviale… (le foncteur $K \mapsto B_K$ ou $K \mapsto \Pi_K$ \uncertain{se factorise} par cette cat. quotient…)}
\note{note écrite en diagonale dans le coin supérieur gauche ; elle concerne vraisemblablement la caractéristique $p$, où l'énoncé de conservativité du texte doit être corrigé des extensions radicielles.}

Quand on prend des corps quelconques, le 2-foncteur $K \mapsto B_K$ ou $K \mapsto \Pi_K$ \add{ou $(K,\overline{K}) \mapsto E_{\overline{K}/K}$} est cependant loin d'être fidèle. Ainsi, si $K$ est sép. clos, $B_K$ est le « topos ponctuel », $\Pi_K$ le groupoïde ponctuel, $E_{\overline{K}/K} \simeq (1)$ -- il est clair que les morphismes entre corps sép. clos ne sont pas décrits par les morphismes entre leurs topos étales, ou groupoïdes fondamentaux ! Pour cette raison, il y a lieu d'associer à un corps $K$ un objet plus fin que $B_K$ ou $\Pi_K$,

\page{6}\note{p. 3 de l'auteur}
à savoir le système projectif des $B_{K_i}$, ou des $\Pi_{K_i}$, pour $K_i$ sous-corps de $K$ de type fini sur le corps premier, et à un syst. $(K, \overline{K})$, le système projectif des $E_{\overline{K}_i/K_i}$, où $\overline{K}_i$ est la clôture alg. sép. de $K_i$ dans $\overline{K}$. On aura
\[
  (\uncertain{4})\qquad
  \begin{cases}
    \Pi_K \simeq \varprojlim \Pi_{K_i} \\
    B_K \simeq \varprojlim B_{K_i} \\
    E_{\overline{K}/K} \xrightarrow{\ \sim\ } \varprojlim E_{\overline{K}_i/K_i}
  \end{cases}
\]
i.e. on reconstitue les objets $B_K$, $\Pi_K$, $E_{\overline{K}/K}$ à partir des syst. proj. correspondants -- mais l'inverse n'est pas vrai.

En fait, comme le foncteur
\[
  \operatorname{Ind}(\text{corps t.f.}) \longrightarrow \text{corps}
\]
de la catégorie des syst. inductifs de corps de type fini, vers celle des corps, est une équivalence de catégories (pour des raisons triviales), il s'ensuit que le foncteur $K \mapsto \mathbb{B}_K$, ou $K \mapsto \boldsymbol{\Pi}_K$, ou $(K, \overline{K}) \mapsto \mathbb{E}_{\overline{K}/K}$, allant
\note{ici les lettres $B$, $\Pi$, $E$ sont tracées doublées ou grasses : elles désignent les systèmes projectifs (pro-objets) qui viennent d'être introduits, par opposition aux $B_K$, $\Pi_K$, $E_{\overline{K}/K}$ ordinaires.}

\page{7}
des corps vers les \ill{} idoines, vont avoir toutes les propriétés de fidélité des foncteurs $K \mapsto B_K$, $K \mapsto \Pi_K$, $(K, \overline{K}) \mapsto \Gamma_{\overline{K}/K}$, restreints aux corps \add{absolument} de type fini, auxquels nous allons par la suite nous borner, le plus souvent du temps. Mais il sera nécessaire, en cours du travail, de donner une description \add{précisément} algébrique, p.ex., des pro-groupes profinis associés p.ex. à $\mathbb{C}$ (plus précisément, à $(\mathbb{C}, \mathbb{C})$ !).

Le rôle dominant sera joué par le corps premier de car. 0, $\mathbb{Q}$, dont le $B_{\mathbb{Q}}$ et $\Pi_{\mathbb{Q}}$, qui a un objet canonique, noté $\overline{\mathbb{Q}}_0$ -- la clôture alg. de $\mathbb{Q}$ dans $\mathbb{C}$. On posera
\[
  (5)\qquad \Gamma_{\mathbb{Q}} = E_{\overline{\mathbb{Q}}_0/\mathbb{Q}}
\]
\marginal{et on écrira souvent $\Gamma_{\overline{\mathbb{Q}}/\mathbb{Q}}$ au lieu de $E_{\overline{\mathbb{Q}}/\mathbb{Q}}$, pour une clôture alg. quelconque $\overline{\mathbb{Q}}$ de $\mathbb{Q}$}
Pour tt corps $K$ de car. 0 -- en particulier pour les corps $K$ de type fini sur $\mathbb{Q}$, aux-

\page{8}\note{p. 4 de l'auteur}
quels nous allons nous borner par la suite -- on a donc des homom. canoniques
\[
  (6)\qquad B_K \longrightarrow B_{\mathbb{Q}}, \qquad \Pi_K \longrightarrow \Pi_{\mathbb{Q}}
\]
qui s'explicitent, quand on a choisi un objet de $\Pi_K$ i.e. un $\overline{K}/K$, disons un $\overline{\mathbb{Q}}/\mathbb{Q}$, par un homom. de groupes profinis
\[
  (7)\qquad E_{\overline{K}/K} \longrightarrow \Gamma_{\overline{\mathbb{Q}}/\mathbb{Q}} .
\]
Par la suite, on regardera toujours $B_K$, $\Pi_K$ et $E$ comme munis de cette structure supplémentaire -- \uncertain{en tant que} les morphismes (de topos, de groupoïdes, de groupes profinis) « arithmétiques » \uncertain{dominant la situation}.

Un intérêt particulier s'attache au noyau de (7), que je note $\pi_{\overline{K}/K}$ -- je l'appelle « partie géométrique » du groupe de Galois $E_{\overline{K}/K}$, par opposition au quotient $E/\pi = \Gamma'_{\overline{K}/K} \hookrightarrow \Gamma_{\overline{\mathbb{Q}}/\mathbb{Q}}$, que j'appelle sa
\marginal{on le notera aussi $\Gamma = \Gamma_{\overline{K}/K}$}

\page{9}
partie « arithmétique » -- celle-ci est un sous-groupe ouvert de $\Gamma_{\overline{\mathbb{Q}}/\mathbb{Q}}$, qui, via la th. de Galois, correspond au sous-corps $k_0$ de $\overline{\mathbb{Q}}/\mathbb{Q}$, extension finie$/\mathbb{Q}$ de $\overline{\mathbb{Q}}/\mathbb{Q}$, clôture alg. de $\mathbb{Q}$ dans $K$, de sorte qu'on a une suite exacte
\marginal{NB $\pi = (1)$ ssi $K$ alg. sur $\mathbb{Q}$, i.e. fini sur $\mathbb{Q}$}

\begin{tikzcd}[column sep=small, row sep=small]
(8)\qquad 1 \arrow[r] & \pi \arrow[r] & E \arrow[r] & \Gamma_{\overline{K}/K} \arrow[r] \arrow[d, hook] & 1 \\
 & & & \Gamma_{\overline{\mathbb{Q}}/\mathbb{Q}} &
\end{tikzcd}

On va donner une interprétation de ce noyau, et de la suite exacte (8), en écrivant
\[
  (9)\qquad K = \varinjlim A_i
\]
où les $A_i$ sont les sous-$\mathbb{Q}$-algèbres de type fini de $K$, correspondant à un syst. projectif des « modèles affines » $U_i = \operatorname{Spec}(A_i)$ de $K/\mathbb{Q}$. Parmi les $A_i$, il y a d'ailleurs un syst. cofinal formé de $A_i$ réguliers, i.e. les $U_i$ lisses$/\mathbb{Q}$, avec comme morphismes de transition des morphismes de localisation $A_i \to A_j = (A_i)_{f_{ij}}$ ($f_{ij} \in A_i \setminus \{0\}$) -- donc les $U_j \hookrightarrow U_i$ des immersions ouvertes, $U_j = (U_i)_{f_{ij}}$. On peut

\page{10}\note{p. 5 de l'auteur}
même, d'après Mike Artin, prendre pour $U_i$ des schémas « élémentaires » sur $k$, se dévissant en fibrations successives de courbes. \add{(Notons que $\operatorname{Spec} K = \eta$ est le pt générique} \struck{Ceci dit, le \ill{}} commun des $U_i$, qui sont géométriquement \add{(est $\Leftrightarrow$ géom. intègres)} sur \struck{\ill{}} $k$ (clôture alg. de $\mathbb{Q}$ dans $K$).
\note{la parenthèse ouverte par l'ajout « Notons que » n'est pas refermée sur la page.}
Le choix de $\overline{K}$ définit un pt géom. $\bar\eta$ sur les $U_i$, d'où des groupes $\pi_1(U_i, \bar\eta) = E_i$, et on a, en vertu de résultats bien connus,
\[
  \operatorname{Spec} K = \varprojlim U_i
\]
\[
  (10)\qquad E_{\overline{K}/K} = \pi_1(\eta, \bar\eta) \xrightarrow{\ \sim\ } \varprojlim_i E_i
  \qquad (\Gamma_i = \pi_1(U_i, \bar\eta))
\]
\note{la parenthèse de droite écrit $\Gamma_i$ là où le texte écrit $E_i$ ; nous laissons l'écart tel quel.}
D'autre part, si on pose
\[
  (11)\qquad \overline{U}_i = U_i \otimes_k \overline{\mathbb{Q}}
\]
on a \add{pour tt $i$} une suite exacte d'homotopie
\[
  (12)\qquad 1 \to \underbrace{\pi_1(\overline{U}_i, \bar\eta)}_{\pi_i}
  \to \underset{\substack{\| \\ E_i}}{\pi_1(U_i, \bar\eta)}
  \to \underset{\substack{\| \\ \Gamma_{\overline{K}/K} = \Gamma_{\overline{\mathbb{Q}}/k}}}{\pi_1(k, \bar\eta)}
  \to 1
\]
\note{le dernier membre de l'égalité sous $\pi_1(k,\bar\eta)$ est surchargé et illisible.}

\begin{tikzcd}[column sep=small, row sep=small]
 & \operatorname{Spec}(K) = \eta \arrow[r] \arrow[d] & \eta_{\overline{\mathbb{Q}}} \arrow[r, no head] \arrow[d, no head] & \bar\eta = \operatorname{Spec}(\overline{K}) \\
 & U_i \arrow[r, no head] \arrow[d, no head] & U_{i,\overline{\mathbb{Q}}} = \overline{U}_i \arrow[d, no head] & \\
\mathbb{Q} \arrow[r, no head] & k \arrow[r, no head] & \overline{\mathbb{Q}} &
\end{tikzcd}

qui forment un syst. projectif de suites exactes, ou d'extensions ayant toutes même quotient $\Gamma'$, et dont les noyaux
\[
  \pi_i = \pi_1(\overline{U}_i, \bar\eta)
\]

\page{11}
sont des groupes fondamentaux « géométriques » -- qui peuvent d'ailleurs se calculer par voie transcendante, en utilisant un plongement de $\overline{K}$ dans $\mathbb{C}$ (d'où un isom $\overline{\mathbb{Q}} \simeq \overline{\mathbb{Q}}_0$), comme les complétés profinis des $\pi_1(U_i(\mathbb{C}), \bar\eta)$, en interprétant maintenant $\bar\eta$ comme un pt \struck{complexe} commun des espaces top. sous-jacents aux variétés complexes $U_i(\mathbb{C})$. La suite exacte (8) est donc la limite projective des suites exactes d'homotopie, ce qui donne en particulier
\[
  \pi_{\overline{K}/K} \simeq \varprojlim_i \underbrace{\pi_1(\overline{U}_i, \bar\eta)}_{\pi_i}
  \simeq \varprojlim_i \pi_1(U_i(\mathbb{C}), \bar\eta)^{\wedge}
\]
\marginal{Cette interprétation \ill{} (12) \ill{} en particulier \ill{} (13) \ill{} ext. \ill{} $K/k_0$}
\note{note en diagonale, à gauche de la formule, très pâle ; seuls les renvois (12), (13) et « $K/k_0$ » se lisent avec quelque sûreté.}
Utilisant les fibrations des $U_i$ (dans le cas où on s'astreint à prendre des variétés élémentaires d'Artin), on trouve que tout $\pi_i$ est un groupe \struck{pr} extension successive de \add{$n$} groupes profinis \emph{libres} (où $n = \dim U_i = \text{deg.tr.}\ K/\mathbb{Q}$). Ceci redonne

\page{12}\note{p. 6 de l'auteur}
p.ex. que la dim. coh. de $\pi_i$ est égale à $n$, celle de $E_i$ est $\leq n+2$ (pour des coeff. \struck{\ill{}} de $m$-torsion, quand $m$ est impair) -- et par passage à la limite, des majorations correspondantes pour les dim. coh. de $\pi_{\overline{K}/K}$ et de $E_{\overline{K}/K}$
\[
  (\uncertain{14})\qquad \operatorname{dim\,coh} \pi_{\overline{K}/K} \leq n, \qquad
  \operatorname{dim\,coh}_\ell E_{\overline{K}/K} \leq n+2 \ \text{si}\ \ell \neq 2
  \quad (\ell\ \text{nb premier})
\]
\note{le numéro ressemble à « (K) » ; nous lisons « (14) » sous réserve, après (12) et (13).}
qui sont en fait même des \emph{égalités} (s.f. erreur), et donnant donc une description cohomologique simple du degré de transcendance absolu de $K$.

\textbf{Théorème 1} Soit $K$ un corps \add{extension} de type fini de $\mathbb{Q}$. \struck{Alors} $\overline{K}$ une clôture algébrique de $K$. Alors pour tout sous-groupe ouvert $E'$ de $E_{\overline{K}/K}$, son centralisateur dans $E_{\overline{K}/K}$ est réduit au groupe unité. Itou pour $\pi_{\overline{K}/K}$.

\emph{Démonstration.} Soit \struck{\ill{}} $\Gamma'$ \add{$\Gamma' \subset \Gamma_{\overline{\mathbb{Q}}/\mathbb{Q}}$} l'image de $E'$ dans $\Gamma$ \add{$= \Gamma_{\overline{K}/K}$}, qui est donc un ss-groupe ouvert. \struck{L'image dans $\Gamma$ du centralisateur de $E'$ \ill{} $E'$ est \ill{}} \struck{si $g \in \Gamma_{\overline{K}/K}$ centralise $E'$, son image $g$ dans \ill{}} \add{\ill{} centralisateur de $\Gamma'$ dans $\Gamma_{\mathbb{Q}}$} \struck{$\Gamma$ centralise $\Gamma'$ \ill{}} Je dis qu'il est
\note{trois lignes biffées et une addition interlinéaire pâle, enchevêtrées ; seul le sens général se lit : on passe à l'image dans $\Gamma$ du centralisateur.}

\page{13}
égale à 1, ce qui équivaut donc au

\textbf{Corollaire} Dans $\Gamma_{\mathbb{Q}} = \Gamma_{\overline{\mathbb{Q}}_0/\mathbb{Q}}$, le centralisateur de tout sous-groupe ouvert est réduit à $(1)$.

OPS le sous-groupe ouvert $\Gamma'$ invariant. Il est bien connu (?) que son centre est réduit à 1 donc si $Z$ est son centralisateur dans $\Gamma_{\mathbb{Q}}$, l'hom $Z \to \Gamma_{\mathbb{Q}}/\Gamma'$ est injectif donc $Z$ est fini. Mais on sait que les seuls éléments $\neq 1$ de $\Gamma_{\mathbb{Q}}$ d'ordre fini sont les conjugués de $\tau$, conjugaison complexe. Mais le centralisateur de $\tau$ dans $\Gamma_{\overline{\mathbb{Q}}}$ est réduit à $(1, \tau)$, donc ne peut contenir $\Gamma'$, donc $\tau \notin Z$, donc $Z = (1)$.
\marginal{? à vérifier}
\note{« OPS » : on peut supposer.}

Revenant à $E' \subset E_{\overline{K}/K}$, on sait donc que son centralisateur $Z$ dans $E$ a une image dans $\Gamma$ réduite à $\{1\}$, donc $Z \subset \pi_{\overline{K}/K}$. Soit $\pi' \subset \pi = \pi_{\overline{K}/K}$ la trace de $E'$ sur $\pi$, c'est un sous-groupe ouvert de $\pi$, et on est ramené à

\page{14}\note{p. 7 de l'auteur}
voir que $\operatorname{Centr}_\pi(\pi') = \{1\}$, i.e. le

\textbf{Corollaire} Soit $\pi$ un groupe profini, extension successive de groupes profinis libres. Alors le centralisateur \add{$Z$ dans $\pi$} de tout sous-groupe ouvert $\pi'$ de $\pi$ est réduit à $\{1\}$.

Par dévissage on est ramené au cas d'un groupe \add{profini} libre. \add{On sait que $\pi'$ est donc libre.} OPS $\pi'$ invariant, et on admet que \struck{?} le centre d'un groupe profini libre est réduit à 1.
\marginal{à vérifier}
Donc $Z \to \pi/\pi'$ est injectif, donc $Z$ est fini, et on admet que dans un groupe profini libre, il n'y a pas d'élément d'ordre fini $\neq 1$
\marginal{à vérifier}
-- ce qui achève la démonstration.

\textbf{Scholie} Le fait que $E_{\overline{K}/K}$ soit \struck{sans} \add{à} \struck{\ill{}} \add{centre} trivial peut s'exploiter en disant que le groupoïde $\Pi_{\overline{K}/K}$ \add{(associé type $B_K$)} est connu à équivalence près, \struck{\ill{}} \add{défini à isom unique près}, quand on connaît le groupe extérieur associé à $E_{\overline{K}/K}$. Les hom

\page{15}
$E_{\overline{K}'/K'} \to E_{\overline{K}/K}$ associés à des hom \add{$K \to K'$} d'extensions de type fini de $\mathbb{Q}$, ayant une image ouverte \uncertain{donc} un centralisateur réduit à 1, on sait de \uncertain{même} que l'hom de topos $B_{K'} \to B_K$ et de groupoïdes $\Pi_{K'} \to \Pi_K$, sont déterminés à \struck{isomorphisme} \ill{} près (définis à isom unique près) par l'hom correspondant des groupes extérieurs. Il en est en particulier ainsi des morphismes structurels $B_K \to B_{\mathbb{Q}}$ et $\Pi_K \to \Pi_{\mathbb{Q}}$, qu'on peut interpréter \add{intrinsèquement} comme un hom de groupes \add{profinis} extérieurs $E_K \to E_{\mathbb{Q}}$. Mais nous préférons suivre un point de vue un peu différent, en exploitant le fait que $\pi_{\overline{K}/K}$ est lui aussi à centre trivial. Cela signifie que l'extension de $\Gamma = \Gamma_{\overline{K}/K}$ par $\pi_{\overline{K}/K}$ est entièrement connue, \add{à isom unique près,} (pour $\pi_{\overline{K}/K}$ \add{et $\Gamma$} fixés)

\page{16}\note{p. 8 de l'auteur}
en termes de \struck{\ill{}} l'action extérieure correspondante de $\Gamma$ sur $\pi$, comme l'image inverse de l'extension universelle

\begin{tikzcd}[column sep=small, row sep=small]
(15)\qquad 1 \arrow[r] & \pi \arrow[r] & \operatorname{Aut}(\pi) \arrow[r] & \operatorname{Autext}(\pi) \arrow[r] & 1 \\
 & & & \Gamma \arrow[u, "\varphi"'] &
\end{tikzcd}

Pour $K$ fixé, donc $k$ fixé, dire qu'on fixe un $\Gamma = \Gamma_{\overline{\mathbb{Q}}/k}$ revient à dire qu'on fixe une clôture alg. de $k$, dire qu'on fixe \add{un} $\pi_{\overline{K}/K} = \pi_1(K \otimes_k \bar k)$ signifie qu'\add{on} fixe un rev. universel de $\operatorname{Spec}(K \otimes_k \bar k) = \eta \otimes_k \bar k$, \struck{compte tenu} les deux ensembles reviennent à se donner le rev. universel \struck{$\operatorname{Spec} K$} $\bar\eta = \operatorname{Spec}(\overline{K})$ de $K$. Par la suite, nous décrivons (avec une fidélité qui reste à examiner) les couples $(K, \overline{K})$ d'une ext. \add{$K$} de $\mathbb{Q}$ de type fini, et d'une clôture algébrique $\overline{K}$ de $K$, par les triples $(\pi, \Gamma, \varphi)$, où $\pi = \pi_{\overline{K}/K}$ et $\Gamma = \Gamma_{\overline{K}/K}$ \add{sont des} groupes profinis, et $\varphi : \Gamma \to \operatorname{Autext}(\pi)$ une action extérieure de $\Gamma$ sur $\pi$ -- ce qui

\page{17}
permet de reconstituer l'extension $E_{\overline{K}/K}$ de $\Gamma_{\overline{K}/K}$ par $\pi_{\overline{K}/K}$. J'ai oublié de préciser qu'il faut de plus se donner $\Gamma$ comme sous-groupe d'un $\Gamma_{\overline{\mathbb{Q}}/\mathbb{Q}}$ bien déterminé, i.e. qu'il faut se donner un objet de $\Pi_{\mathbb{Q}}$ et une action fidèle de $\Gamma$ dessus -- pour reconstruire par l'ensemble de ces données un hom de groupoïdes profinis $\Pi_K \to \Pi_{\mathbb{Q}}$, plus un objet de $\Pi_K$ -- ou encore, un morphisme de topos pro-galoisiens $B_K \to B_{\mathbb{Q}}$, plus un pt de $B_K$.

On peut aussi fixer un objet de $\Pi_{\mathbb{Q}}$, i.e. un pt de $B_{\mathbb{Q}}$, i.e. un $\overline{\mathbb{Q}}$, et étudier les $K$, \struck{tels que $k$} \add{avec} un plongement de $k$ (clôture alg. de $\mathbb{Q}$ dans $K$) dans $\overline{\mathbb{Q}}$ -- \struck{ils sont} mais sans se donner une clôture alg. $\overline{K}$ de $K$ qui induise $\overline{\mathbb{Q}}/\mathbb{Q}$. Ils sont décrits

\page{18}\note{p. 9 de l'auteur}
On a ainsi plusieurs points de vue essentiellement équivalents, pour décrire par voie profinie \struck{une} \add{$K$} extension de type fini de $\mathbb{Q}$ :
\begin{enumerate}
  \item Par le topos étale $B_K$, en tant que topos pro-galoisien sur $B_{\mathbb{Q}}$.
  \item Par le groupoïde fondamental $\Pi_K$ de ce topos (le groupoïde de ses pts, ou de ses rev. universels) -- en tant que groupoïde au-dessus de $\Pi_{\mathbb{Q}}$.
  \item Par le groupe extérieur $E_K$, au-dessus du groupe extérieur $E_{\mathbb{Q}} \simeq \Gamma_{\mathbb{Q}}$ (pro-groupe).
  \item En termes d'une clôture algébrique $\overline{K}/K$ (i.e. en décrivant des couples $(K, \overline{K})$ plutôt que $K$), par un objet $\overline{\mathbb{Q}} \in \operatorname{Ob} \Pi_{\mathbb{Q}}$, et un hom de groupes profinis $E \to \Gamma_{\overline{\mathbb{Q}}/\mathbb{Q}}$. \struck{un sous-groupe ouvert $\Gamma$ de}
  \item En termes d'une clôture alg. \add{fixe} $\overline{\mathbb{Q}}$ de $\mathbb{Q}$, et si $\Gamma = \Gamma_{\overline{\mathbb{Q}}/\mathbb{Q}}$, \struck{par le groupe} en considérant les couples $(K, i)$ où $i : k \to \overline{\mathbb{Q}}$ est un plongement de la clôture alg. \add{$k$} de $\mathbb{Q}$ dans $K$ dans $\overline{\mathbb{Q}}$ : \struck{\ill{}} par le groupe extérieur $\pi_K = \pi_1(K)$, sur lequel un sous-groupe ouvert $\Gamma_K \subset \Gamma$ opère extérieurement.
\end{enumerate}
\note{la numérotation est de l'auteur, de 1) à 5) ; un premier « 2) » est biffé.}
\marginal{6) En termes d'un $\overline{\mathbb{Q}}/\mathbb{Q}$ : par le groupoïde $\Pi_{K \otimes_{\mathbb{Q}} \overline{\mathbb{Q}}}$ (\ill{}) sur lequel $\Gamma_{\overline{\mathbb{Q}}/\mathbb{Q}}$ opère -- ou \ill{} le topos galoisien $B_{K \otimes_{\mathbb{Q}} \overline{\mathbb{Q}}}$, \ill{} profinis, \ill{} $\Gamma_{\overline{\mathbb{Q}}/\mathbb{Q}}$ opérant dessus…}
\note{note verticale dans la marge gauche, en partie illisible.}

\page{19}
par des groupes \add{profinis} extérieurs $\pi_1(K) = \Gamma_K$, sur lesquels \add{un sous-groupe ouvert $\Gamma$ (non précisé d'avance) de} $\Gamma_{\overline{\mathbb{Q}}/\mathbb{Q}}$ opère extérieurement.
\note{la phrase reprend la fin de la p. 17 (« Ils sont décrits »). Le $\Gamma_K$ de « $\pi_1(K) = \Gamma_K$ » est tel sur la page, là où l'on attendrait $\pi_K$.}

Un hom de corps $K \to K'$ donne lieu à un hom de groupes extérieurs, \struck{\ill{}} $\pi' \to \pi$, où l'image de $\pi'$ dans $\pi$ est ouverte donc de centralisateur réduit à $(1)$, ce qui \struck{signifie} \add{implique} encore que le morphisme de topos $B_{K' \otimes_k \overline{\mathbb{Q}}} \to B_{K \otimes_k \overline{\mathbb{Q}}}$ est déterminé (à isom unique près) par cet hom extérieur. De plus on a des actions extérieures de $\Gamma = \Gamma_K \subset \Gamma_{K'}$ \add{sur} $\pi'$ et $\pi$, de façon que $\pi' \to \pi$ y commute -- et ceci suffit pour reconstituer, d'une part les groupes extérieurs $E$, $E'$ extensions \add{(« extérieures »)} de $\Gamma$ par $\pi$, $\pi'$ (et, à équivalence rigide près, les $B_K$, $B_{K'}$ et $B_K \to B_{\mathbb{Q}}$, $B_{K'} \to B_{\mathbb{Q}}$) et de plus l'hom d'extensions \add{(extérieures)} $E \to E'$ de $\Gamma$.
\marginal{\ill{} $k \subset \overline{\mathbb{Q}}$ \ill{}}

\page{20}\note{p. 10 de l'auteur}
\textbf{Remarques} Quand $\pi \neq (1)$, \add{i.e. $K$ pas fini sur $\mathbb{Q}$,} le théorème 1 peut se renforcer, \add{sauf erreur,} en écrivant que pour tout sous-groupe \add{$\pi' \subset \pi$} ouvert \emph{dans $\pi$}, $\operatorname{Centr}_E(\pi') = \{1\}$. Si $Z$ est ce centralisateur, on a $Z \cap \pi = (1)$ d'après le Th. 1, \struck{et} \add{prouvons que} l'image de $Z$ dans $\Gamma_{\overline{K}/K} \subset \Gamma_{\overline{\mathbb{Q}}/\mathbb{Q}}$ est \struck{contenue dans} \add{finie} \struck{\ill{}} (ce qui prouvera \add{déjà}, \uncertain{dans}, que $Z$ est d'ordre 1 ou 2, et dans le deuxième cas que son image dans $\Gamma_{\overline{\mathbb{Q}}/\mathbb{Q}}$ est \struck{\ill{}} engendrée par un $\tau$ de conjugaison complexe).

Quitte à remplacer $E$ par un sous-groupe ouvert assez petit (ce qui revient à passer à une extension finie de $K$) OPS $\pi' = \pi$, alors l'image $Z'$ de $Z$ dans $\Gamma$ est contenue dans le noyau de l'hom.
\[
  \varphi : \Gamma \longrightarrow \operatorname{Autext}(\pi) .
\]
Or (sauf erreur) je sais prouver que cet hom. est injectif (on est ramené aussitôt au cas où $K$ est de degré de transc. 1, et on est ramené au cas \add{du $\pi_1$} d'une courbe algébrique…)
\note{la page s'achève ici ; la remarque se poursuit peut-être au-delà du lot.}

\end{document}
